\chapter{Can We Think Without Words?}
\section{The Hand Axe Behind Glass}
First, pick something up.

Not a phone or a book, but a stone. You may have seen it in a natural-history museum: teardrop-shaped, slightly longer than an adult palm, neatly flaked around the edge, almost symmetrical. The label says hand axe, and its age is startling: hundreds of thousands, perhaps more than a million years.

Look for ten more seconds and notice the symmetry. Falling or trampling did not make this stone. Its maker took flint and struck away dozens of flakes at different angles. Before each blow, the maker had to compare the present shape with a mental target, a symmetrical teardrop not yet in existence. A wrong blow could ruin it irreversibly. Before the stone became an axe, an axe's shape occupied a mind, perhaps for hours, repeatedly compared, revised, and approached.

That mind belonged to an ancient human hundreds of thousands of years ago. Did its owner have language like ours? Nobody knows. Language leaves no fossils, and vocal cords decay far faster than bones. Many researchers suspect that even if those humans communicated, they lacked a full language for yesterday, tomorrow, and if.

What, then, was the teardrop target in that craftsperson's mind through hours beside the stone? If it was thought, thought occurred outside words. If not, we need another name, yet it plainly involved planning, imagination, correction, and patience, things we ordinarily count as thinking.

Now the other extreme. A myna in a market cage calls hello to you. It has a word. But does it have the intention hello points to? Is it greeting you or playing a recording?

One perhaps has no words but holds a plan; the other plainly has words but seems to run empty. Between them lies our question: \textbf{Can people think without words? Conversely, does having words guarantee thinking?}

\section{A Classroom Without Sound}
Now enter a scene.

Time: the late 1970s. Place: a new school in Managua, Nicaragua's capital. The country has just emerged from civil war and fresh slogans cover walls. For the first time, the government opens schools for deaf children. Dozens, then hundreds arrive from across the country and sit in the same classrooms.

The classroom looks unusual. Teachers write Spanish words on the board, exaggerate mouth movements, and repeatedly point to their lips. Their task is to teach lip-reading and spoken Spanish. Children watch politely and learn little. This resembles their years at home: born into hearing families, surrounded by conversations in a language they cannot hear, living as if under glass.

But move your attention from the teacher to the playground and school buses. There, hands fly through the air, fingers, palms, eyebrows, shoulders moving like a downpour. The children are chatting enthusiastically.

An unprecedented situation has emerged. Most children already invented home signs, a few dozen gestures worked out with parents for eat, want, where. That itself is part of an answer. Without anyone teaching them a language, they urgently wanted to express themselves and invented means. Now hundreds of home-sign systems meet on one playground. Older children borrow and combine them into shared signs, lively but rough, like a temporary plank walkway.

What astonishes linguists is the younger cohort arriving afterward. They learn the shared system, but their signing differs from what older children supplied. Movements divide into smaller components, components acquire fixed orders, and a sign takes different roles in different sentences. \textbf{Grammar emerges.} Teachers did not teach it; they do not understand the system. Older children did not have it. Younger children develop it through use. Over subsequent decades, linguists such as Ann Senghas follow the school, watching the language grow and become more complex with each cohort, like a seed sprouting under a microscope. It is now called Nicaraguan Sign Language.

Hold this scene in memory. Children without a ready-made language first think, want, and invent individually, then together create a full grammatical language within decades. Thought runs ahead of words, then draws them into being.

\section{Three Questions Collide}
Put the conflict onstage without deciding it yet.

We have three seemingly conflicting things. A hand-axe maker perhaps without full language holds a goal requiring planning, imagination, and correction. A myna holds words without seeming to hold thoughts. Deaf children in Managua, taught no language by anyone, first think and express themselves, then create one.

They illuminate three sides of a question that is really several questions:

\begin{itemize}
\item \textbf{Does thought need words?} Is the incessant inner voice, what shall I eat, did I say that wrong, thought itself or merely its clothing?
\item \textbf{Where did language come from?} An accident of sound or a legacy of gesture? A gift switched on by mutation or a tool refined by generations of children?
\item \textbf{Why have only humans turned communication into language?} Bees send messages, monkeys alarms, dogs understand hundreds of words. Which components separate them from us?
\end{itemize}
These questions intertwine; we will separate them. First establish a rule for the whole book: distinguish what happened, evidenced facts; why, competing explanations; how participants understood it, their accounts; and how we evaluate it, values. Confusing these produces much unnecessary argument.

Begin with the first question: do words live inside thought, or thought inside words?

\section{Both Sides Have Reasons}
There are two serious answers to whether thinking needs words. Before any conclusion, something more valuable: both deserve their strongest case before comparison.

\textbf{Position one: without language, no substantial thought.}

Let me present this fully. Our age often dismisses it with language is only a tool, but its cards are strong.

First, your own experience. Mentally multiply thirty-seven by six. Did you notice an inner voice reciting six sevens are forty-two, write two and carry four, six threes are eighteen? Remembering the rule and tracking the carry depend on internal language. Without number words and this verbal procedure, an ordinary person cannot do the calculation. How do you remember a phone number? Repeat it internally: language again. Much complex thinking genuinely is talking inside our heads.

Second, some things are difficult to grasp without words. Fairness, inflation, empathy, involution: not visible objects, but bundles of vague feelings tied together and labeled. The label is a word. Without it, the bundle scatters. You cannot pass it to others, or even reliably return it to yourself days later. Words are thought's handles.

Third, civilization. However clever a chimpanzee, its intelligence disappears with its death. Humans store ideas in words, then books and networks, so each generation begins within the previous generation's mind. Without language as external storage, knowledge cannot accumulate.

\textbf{Position two: thought comes first; words merely carry it.}

This side has strong evidence too.

Long before a first word, babies already think. An eight-month-old lifts a cloth covering a toy, knowing the unseen toy remains. Object permanence is genuine reasoning, appearing months before the word ball.

You encounter counterexamples daily. A thought reaches your tongue but words will not come. That moment shows the idea is already present: you know what you mean and can recognize whether someone else expresses it correctly, but the words have not arrived. Meaning comes first. Or a geometry problem suddenly yields the auxiliary line you need after long staring. First you see it, then find a description. In the instant of insight, no sentence is present.

Stronger evidence comes from hospital wards. Neurologists have long observed people whose brain injuries remove language. They cannot form sentences or understand complex ones, yet many can play chess, calculate, navigate, plan a day, and recognize deception. The language room has collapsed, but its occupant remains. Recent brain-imaging studies point similarly: regions handling language differ from those active in arithmetic and logical reasoning; language areas can be quiet during logic problems.

Finally, remember Managua. If thought had to wait for words, the deaf children would have been empty shells. Instead, their long-unexpressed ideas forced a language into existence. Language did not give birth to thought; thought gave birth to language.

An honest formulation is therefore: \textbf{thinking is not identical to inner speech, but inner speech is an especially useful kind of thinking}. Language resembles a bicycle. Legs already walk, but wheels take you farther. We will return to this image.

\textbf{Another set of voices: what did language first look like?}

If deaf children can create a full language with their hands, the assumption that language equals speech must go. Language's core is symbols and grammar; sound is merely its most common carrier. Two major origin hypotheses follow. These are hypotheses throughout: language leaves no fossils, and nobody has footage.

The gesture camp argues that monkeys and apes have limited control over vocal production but good hand control. Captive apes can learn hundreds of signs. Perhaps ancestors first spoke with hands and later changed channels to the throat. Managua shows that pure gesture can become a complete language. The vocal camp replies that hands work while mouths are free; talking while gathering fruit is socially efficient. Sound also travels in darkness or behind visual barriers, and human throats have specialized changes for vocalization. Both may be partly right: gesture begins, sound takes over. Neither can produce fossils to decide it. This is the chapter's thinnest evidence, but marking what we do not know is part of humanistic training.

\section{Thinking Bridge: Three Thresholds}
My dog understands everything and animals have no language collide in ordinary arguments because the speakers use different measures. Here is a portable ruler. Comparing human languages with known animal communication, linguists identify key features. We turn three especially useful ones into thresholds. For any system, animal calls, bee dances, memes, alien signals, do not ask merely whether it means something. Ask whether it passes these three gates.

\textbf{First, combination.} Finite components form indefinitely many new sentences according to rules. Most sentences you hear are new: yesterday I saw a cat in a raincoat in the elevator. You understand immediately because meaning resides in components and their arrangement. Dog chases cat and cat chases dog use the same pieces but reverse the meaning. Parts, order, and construction jointly determine meaning. This is language's extraordinary efficiency: a few thousand common words serve a lifetime of new sentences.

\textbf{Second, openness.} Can the system address anything? Yesterday's new word, next year's plan, unseen animals, nonexistent dragons, unicorns, extraterrestrial deliveries? Human language puts no limit on topics. New things promptly receive words, and old words acquire new uses.

\textbf{Third, displacement.} Can it address things absent from here and now? This easily overlooked feature matters enormously. Yesterday's leopard, if it rains tomorrow, he said he was sorry though he was not: past, future, hypothesis, fiction, and lies all require absent objects. Language releases us from the prison of the present situation.

Now measure animals against these gates.

Vervet monkeys on African grasslands are stars of communication research. Different calls warn of leopards, eagles, and snakes. Companions climb trees at the leopard alarm and stand to inspect grass at the snake alarm. Recorded playback produces the same responses, showing that information lies in the call rather than the scene's atmosphere. Impressive. Yet the inventory is closed, just a few calls without combination, novel messages, or conversations about yesterday's leopard. Full marks for communication, but none of the three gates passed.

An easily misjudged counterexample comes next. Textbooks say only humans discuss absent things. Bees disagree. A worker returning from nectar performs a waggle dance whose orientation indicates direction relative to the sun and whose waggle duration indicates distance. Food kilometers away is absent. This is genuine displacement. But the other gates stay shut: the dance always concerns food over there. It cannot discuss yesterday's deceptive flower or suggest canceling a trip if rain comes. One wide-open gate does not make three.

Dogs? Experiments showed that Rico, a border collie, understood more than two hundred object names and could infer a new word. Hearing an unfamiliar name, he selected an unfamiliar toy from familiar ones, using exclusion to guess its meaning. Exciting, but Rico could not speak a word. Understanding two hundred nouns and combining components into unlimited sentences are different abilities. Listening takes him halfway through a gate; speaking and combining offer no gate at all.

Alex the parrot went farther. After decades of training, he recognized colors and shapes, counted small quantities, and combined words into simple requests. Yet honesty is needed: decades of deliberate training yielded a fraction of an ordinary three-year-old's vocabulary and grammar, while that child had no comparable special training and developed it naturally.

Closest are apes learning signs or symbols from humans. Some learn hundreds of signs, others request food on symbol boards, occasionally producing novel two-word combinations. But the history of checking these claims is equally famous. Reviewing footage from a sentence-teaching project revealed that many apparent utterances imitated a trainer's preceding actions or responded to unconscious cues, like Clever Hans, the supposedly calculating horse who read spectators' expressions. Debate remains, but even optimistic readings leave these systems at dozens or hundreds of pieces and very simple combinations, far from unrestricted expression.

The conclusion has two halves. So far only humans pass all three gates, including those using hands: Managua reminds us language is not sound. Equally important, \textbf{no language never means no mind}. Monkey alarms report the world, dogs' exclusion is reasoning, bees' dances are displacement. They are thinkers, simply not thinkers who express thought as we do. Treating animals as mindless machines is as lazy as treating history as predetermined.

\section{Two Statements from 2,300 Years Ago}
The question is not new. A primary source shows how deeply Warring States thinkers explored it.

In Xunzi's ``Correct Naming,'' a passage asks where names come from:

\begin{quote}
Names have no fixed appropriateness. They are assigned by agreement; established convention makes them appropriate, and departure from that agreement makes them inappropriate.
\end{quote}
Roughly, no name naturally fits. People agree that a sound will indicate a thing; repeated use becomes custom and establishes the name. Using it outside the agreement does not.

Pause over its sharpness. Why is a dog called gou in Chinese? Not because that character is inscribed on the animal or the sound resembles it. English uses dog, Japanese inu. Different agreements all work. No natural cord joins word to thing, only a community's convention. Today this is called the arbitrariness of the sign, a foundation of linguistics that Xunzi articulated 2,300 years ago.

He goes deeper. If convention makes a name appropriate, who participates, who may name, and whether old agreements should yield to new usage become social questions. Meanings belong not to lexicographers but to millions of users. Thus neijuan, involution or exhausting internal competition, and shekong, social anxiety, can establish themselves without approval. Agreements are renewed daily.

Now another sentence, from the Zhuangzi's ``External Things'':

\begin{quote}
Words are for grasping meaning; once meaning is grasped, words may be forgotten.
\end{quote}
Language catches meaning; after catching it, you can put the words down. The original uses fishing and rabbit-catching images: a trap catches fish, and once the fish is caught the trap may be forgotten. Do not cling to words after understanding.

Together, Xunzi and Zhuangzi each manage half this chapter. Xunzi says words are agreed tools and agreement is social, answering what language is. Zhuangzi says meaning is the quarry and words the trap, answering which is larger, language or thought. Ancient Chinese thinkers distinguished words from meanings, in striking accord with evidence from wards and playgrounds: the quarry precedes the trap.

We cite them not because everything ancient thinkers said was right. Many of their ideas about language's origins do not hold today. They supply a measure of time: struggling with your inner voice means joining a meeting at least 2,300 years old.

\section{How Do Babies Learn Without Grammar Class?}
Now establish a central fact and offer different explanations.

The factual starting point is presented as undisputed: babies in Beijing, Nairobi, and Lima follow remarkably similar language timetables. Around six months, babbling tunes their vocal apparatus; around one year, the first word; between eighteen months and two years, two-word expressions such as Mommy up or dog go. Vocabulary and sentence complexity then surge, and by four or five children fluently produce and understand \textbf{sentences never heard before}. They are sentence-makers, not recorders. A girl says she spilled the milk because the table bumped it. Nobody taught her that sentence. A preposition may be unconventional, but she generated the grammatical logic herself.

No grammar lesson accompanies the process. In some cultures, adults do not even address infants especially; children learn while immersed in surrounding conversation. Teaching cannot explain it because there is no classroom.

Why? Here are three genuinely different explanations.

\textbf{Explanation one: a language device is preinstalled in the brain.}

In the mid-twentieth century, Noam Chomsky argued that children's input is limited and messy, full of errors, fragments, and indistinct pronunciation, yet their resulting grammars are highly consistent. If input cannot explain output, the remainder must come from the child: a shared deep framework already present in human brains, with learning setting switches to Chinese or English. A strong piece of evidence for universal grammar is the critical-period phenomenon, as if language has an expiring ticket. Managua illustrates it. Earlier exposure yields fuller grammar; deaf adults first encountering a language can acquire vocabulary yet face grammatical ceilings. Genie, isolated until rescue at thirteen, later learned many words but never complete syntax. Her suffering requires restraint in interpretation: she lacked not merely language input but a whole life. Her case supports, but cannot decide, the argument.

The limit is specifying the preinstalled framework. Lists of universal features have shortened through decades of revision, with no agreed complete version. Innate can also become a shield against explanation. Attribute everything unexplained to nature and science stops.

\textbf{Explanation two: children are statisticians and psychologists.}

Other researchers focus on learning skills and social instincts. They include developmental psychologist Michael Tomasello. Statistical-learning experiments supply evidence. Eight-month-olds hear two minutes of meaningless syllables, without pauses or stress cues; only certain syllables consistently occur together. Tests show they distinguish frequent combinations, possible words, from random ones. Two minutes, eight months old, no teacher. Babies begin segmenting words by statistics, seemingly without a dictionary. This account also sees them as mind-readers, following adults' gaze and pointing to infer intentions and connect sounds to meanings. Language is not a device descending from above but patterns extracted from countless exchanges of pointing and looking.

Its limit: statistics can form apparently permissible sentences, but children also reject unheard sentences that feel wrong. The origin of this intuition for abstract rules remains a challenge the account continues to address.

\textbf{Explanation three: children have filtered language.}

This reverses the causal question. Instead of asking why children's brains fit language, ask why language fits children's brains. Unfitting forms disappeared. Each transmission passes through children; hard structures become simplified, regularized, and gradually eliminated. Language is not a tailored suit for the brain but old jeans worn into shape by generations. Managua again supplies living evidence: successive younger cohorts turn older children's plank walkway into a paved road. Every generation of learners edits language.

Its limit: it explains what language \textbf{becomes}, not who began the process. Why human ancestors started it still returns us to evolution's puzzle.

These are not mutually exclusive choices. Brains may have innate learning biases, statistics and social interaction may build language daily, and generations may repeatedly reshape the language itself. Distinguishing the accounts is not about choosing a camp, but seeing how one agreed fact supports several serious disputes. Facts and explanations occupy separate levels. Someone who recognizes them is harder to fool in any debate.

\section{Further Challenge: Tracks for Thought}
Now a famous linguistic theory, important enough to have a familiar shorthand and tempting enough to be misused.

In the first half of the twentieth century, Edward Sapir and his student Benjamin Lee Whorf advanced a large question: do speakers of different languages see different worlds? Linguistic relativity, often called the Sapir--Whorf hypothesis, has two versions of very different strength. Separating them is essential.

\textbf{The strong version: language determines thought.} You can think only what your language can express. Evidence has largely defeated this. If true, Managua's children would have had no thoughts before signing, yet their thoughts helped create the language. Tip-of-the-tongue states should not exist, since thoughts would always arrive with words. It also implies cultural prisons, each language community trapped in words and unable to understand others. Translation, difficult but real, contradicts that, and it makes cultures into sealed boxes, a recurring danger in this book.

\textbf{The weak version: language influences habits of thought.} It builds paths, not walls. Distinctions your language regularly requires become quicker, more frequent, and less effortful. Three elegant kinds of experimental evidence illustrate this.

First, color. Russian has no single basic blue term: siniy for dark blue and goluboy for light blue have separate basic status, like red and yellow in Chinese. Experiments find faster Russian than English discrimination between blues lying on opposite sides of that boundary, despite the same eye structure. A lexical boundary becomes a perceptual speed bump, measured in milliseconds and operating daily.

Second, direction. Guugu Yimithirr is an Indigenous language of Queensland; its gangurru gave English kangaroo. It largely uses compass directions instead of left and right. Speakers say roughly the cup north of you or watch the ant south of you. Walking, talking, and storytelling therefore require an internal compass. Tests find striking directional accuracy even in unfamiliar places or after blindfolded turning. This is not extraordinary natural talent; their language supplies constant practice.

Third, number. The Amazonian Munduruku have number words extending approximately to five, then some and many. Researchers find intact approximate quantity comparison, such as choosing the larger pile of beans, but difficulty beyond five with \textbf{exact} calculations such as seven minus four. This offers a precise answer to our opening question. Without larger number words, what is lost is not quantity sense but \textbf{exactness}. Approximate amounts can be thought clearly; exact numbers depend on words. Language here is a ruler, not an on-off switch.

Notice an \textbf{easily misjudged counterexample}: a famous argument for relativity famously failed. Whorf claimed Hopi lacked time words and tense, and that Hopi people therefore lacked our concept of time. The claim spread widely and still appears in popular books. Later, the linguist Ekkehart Malotki actually learned Hopi and wrote a substantial book documenting abundant time vocabulary, tense marking, and temporal expressions. The supposedly timeless people had perfectly serviceable schedules. The strong version's famous witness did not hold up.

Remember this failure because it shows a seductive error: no word for something means no ability to think it. By that logic, English lacking a single exact equivalent of Chinese xiao, filial piety, would make British people unable to understand caring for parents. Absurd. Words are convenient paths, not walls. They affect speed, not possible destinations. Yet paths matter: build one and traffic increases. Once neijuan names a pervasive, hard-to-describe exhaustion, people can discuss and compare it. The word need not invent the feeling, but makes it public. Language builds highways, not prisons, and their placement quietly alters traffic.

Weigh the boundary. The theory explains efficiency differences in color, direction, and number well, who performs the same task a little faster. It cannot settle grand claims about whether thought is possible, much less declare how a people naturally thinks. Cultural determinism uses a precise tool for crude work.

\section{When the Inner Voice Disappears}
Today, this ancient meeting continues in your mind and on your phone.

First, a difference many never notice: inner speech varies enormously. Recent psychological surveys find some people with all-day narration, even a voiceover while reading, and others whose inner lives are largely quiet, relying on images, feelings, or direct knowing. Both think, take exams, fall in love, and repair bicycles normally. This is the chapter's conclusion in everyday form: inner speech is a way of thinking, not a thinking license.

Small evidence hides in daily experience. At the tip of your tongue, meaning is present and words absent. Mental arithmetic borrows words, but suddenly seeing an auxiliary line does not. The painful essay-writing moment, I have it inside but cannot write it, separates thought and words. The quarry is in hand before the trap is woven.

Now the new-word production line. Recent Chinese expressions such as shekong, social anxiety; qingxu jiazhi, emotional value; and zuiti, someone who voices what you want to say, do what Xunzi described: agreement. Millions vote with usage, binding sounds to feelings until convention establishes words. No dictionary committee's permission is required. You witness language becoming conventional in a natural experiment ancient thinkers could only imagine.

The internet also poses another question: are memes a language? Apply the gates. Some combination: images can form new meanings. Displacement: memes often concern absent things. Openness falls short: they cannot readily discuss a second-quarter budget or specify what follows if an apology is rejected and fault is limited. It is real, clever communication, but incomplete, rich enough for feeling and poor for contracts. Measuring replaces an endless yes-or-no argument with a clearer account of the system.

The newest variable is artificial intelligence. Today's large language models are presented here as pure word machines, taking words in and sending words out. Keep three levels separate: the stated factual claim is that they process language; a reasonable dispute concerns understanding, over which researchers disagree strongly; an open question asks how far a merely speaking machine is from thought if thought is not identical to inner speech. This chapter does not decide. But you have the three gates, quarry and trap, and the distinction between strong and weak relativity, enough to ask better questions than many spectators.

\section{Over to You: Does an Unnamed Feeling Exist?}
Return to the glass case.

The craftsperson left no sentence. Dead for hundreds of thousands of years, perhaps without surviving bones, they left a symmetrical teardrop. Across that distance, through glass, you see the not-yet-an-axe that once lived in their mind. Is it an unspoken sentence?

Now answer a larger version. Think first, then give reasons:

\textbf{If no word yet exists for a feeling, does the feeling exist?}

Make the yes case. Babies feel hurt before speech; Rico wanted to play before understanding ball; Munduruku speakers can sense seven exceeds five without exactly counting seven; tip-of-the-tongue experience puts meaning before words. Feelings are quarry running in the forest whether traps exist or not.

Make the not-necessarily case too. Before neijuan, that exhaustion was scattered, private, inexpressible, each person feeling alone. The word made it one thing people could identify and discuss together. Words summon public feelings. The quarry always ran, but only a finished trap let humanity catch it. For a society, does an uncaught feeling count as existing?

Add a question you may wish to avoid. If wordless feelings count, do a dog's, pig's, or octopus's pain and joy count as thought when they have no words to report them? Should our treatment change? Whichever answer you choose, explain why.

Do words discover feelings or invent them? There is no standard answer. But next time something stirs inside and no word can contain it, you will know where you stand: at the border of thought and language, holding a map not yet printed. Humanity waited hundreds of thousands of years for someone willing to name such a feeling. Perhaps the next word is yours to make.
